% =====================================================================
%  StockIT — Rapport de présentation détaillé
%  Compilation : pdflatex rapport_presentation.tex (x2 pour la ToC)
%  Testé sur Overleaf (pdfLaTeX) et TeX Live.
% =====================================================================
\documentclass[11pt,a4paper]{article}

% ---- Encodage & langue ---------------------------------------------
\usepackage[utf8]{inputenc}
\usepackage[T1]{fontenc}
\usepackage[french]{babel}
\usepackage{lmodern}
\usepackage{microtype}

% ---- Mise en page ---------------------------------------------------
\usepackage[a4paper,margin=2.2cm,top=2cm,bottom=2cm]{geometry}
\usepackage{parskip}
\setlength{\parskip}{4pt}
\setlength{\parindent}{0pt}

% ---- Extensions -----------------------------------------------------
\usepackage{xcolor}
\definecolor{primary}{HTML}{1F4E79}
\definecolor{accent}{HTML}{C0504D}
\definecolor{soft}{HTML}{E8EEF7}
\definecolor{softgrey}{HTML}{F5F5F5}
\definecolor{green1}{HTML}{2E7D32}
\definecolor{orange1}{HTML}{E65100}

\usepackage{enumitem}
\setlist[itemize]{leftmargin=1.3em,itemsep=2pt,topsep=3pt}
\setlist[enumerate]{leftmargin=1.5em,itemsep=2pt,topsep=3pt}

\usepackage{tcolorbox}
\tcbuselibrary{skins,breakable}

\usepackage{tikz}
\usetikzlibrary{arrows.meta,positioning,shapes.geometric,shadows,fit,backgrounds}

\usepackage{booktabs}
\usepackage{array}
\usepackage{tabularx}
\newcolumntype{L}[1]{>{\raggedright\arraybackslash}p{#1}}

\usepackage{titlesec}
\titleformat{\section}
  {\normalfont\Large\bfseries\color{primary}}
  {\thesection.}{0.6em}{}
\titleformat{\subsection}
  {\normalfont\large\bfseries\color{primary!85!black}}
  {\thesubsection}{0.6em}{}
\titleformat{\subsubsection}
  {\normalfont\normalsize\bfseries\color{accent}}
  {}{0em}{}
\titlespacing*{\section}{0pt}{14pt}{6pt}
\titlespacing*{\subsection}{0pt}{10pt}{4pt}
\titlespacing*{\subsubsection}{0pt}{6pt}{2pt}

\usepackage{fancyhdr}
\pagestyle{fancy}
\fancyhf{}
\renewcommand{\headrulewidth}{0.3pt}
\renewcommand{\footrulewidth}{0pt}
\fancyhead[L]{\small\color{primary!80!black}\textbf{StockIT} — Rapport de présentation}
\fancyhead[R]{\small\color{black!60}\thepage}

\usepackage[hidelinks,colorlinks=true,linkcolor=primary,urlcolor=primary]{hyperref}

% ---- Styles TikZ ----------------------------------------------------
\tikzset{
  box/.style   = {rectangle, rounded corners=4pt, draw=primary,
                  fill=soft, minimum height=11mm, minimum width=26mm,
                  align=center, font=\small},
  cbox/.style  = {rectangle, rounded corners=4pt, draw=orange1,
                  fill=orange!10, minimum height=11mm, minimum width=30mm,
                  align=center, font=\small},
  aibox/.style = {rectangle, rounded corners=4pt, draw=accent,
                  fill=red!5, minimum height=11mm, minimum width=30mm,
                  align=center, font=\small},
  layer/.style = {rectangle, rounded corners=3pt, draw=primary!70,
                  fill=soft, minimum height=8mm, minimum width=42mm,
                  align=center, font=\small},
  arr/.style   = {-{Latex[length=2.5mm]}, thick, draw=black!55}
}

% ---- Boîtes réutilisables ------------------------------------------
\newtcolorbox{infobox}[1][]{
  enhanced, colback=soft, colframe=primary,
  boxrule=0.4pt, arc=3pt, left=8pt, right=8pt, top=5pt, bottom=5pt,
  fonttitle=\bfseries\color{white},
  coltitle=white, colbacktitle=primary,
  title=#1, breakable}

\newtcolorbox{warnbox}[1][]{
  enhanced, colback=orange!5, colframe=orange1,
  boxrule=0.4pt, arc=3pt, left=8pt, right=8pt, top=5pt, bottom=5pt,
  fonttitle=\bfseries, coltitle=white, colbacktitle=orange1,
  title=#1, breakable}

\newtcolorbox{ideabox}[1][]{
  enhanced, colback=green!4, colframe=green1,
  boxrule=0.4pt, arc=3pt, left=8pt, right=8pt, top=5pt, bottom=5pt,
  fonttitle=\bfseries, coltitle=white, colbacktitle=green1,
  title=#1, breakable}

% =====================================================================
\begin{document}

% ---------- Page de garde -------------------------------------------
\thispagestyle{empty}
\begin{center}
\vspace*{1.5cm}
{\color{primary}\rule{\linewidth}{1.2pt}}\\[6pt]
{\Huge\bfseries\color{primary} StockIT}\\[6pt]
{\Large\itshape Application mobile intelligente de gestion\\
de stock IT, assistée par IA}\\[4pt]
{\color{primary}\rule{\linewidth}{1.2pt}}

\vspace{2.5cm}

\begin{tcolorbox}[colback=soft, colframe=primary,
  boxrule=0.5pt, arc=5pt, width=0.85\linewidth, halign=center]
\Large\bfseries Rapport de présentation détaillé\\[4pt]
\normalsize\normalfont
Architecture, fonctionnalités, workflows,\\
automatisations et intégrations
\end{tcolorbox}

\vspace{3cm}

\begin{tabular}{rl}
\textbf{Plateforme}      & Android (Java, min SDK~26 — Android~8)\\[3pt]
\textbf{Domaine}         & Logistique / Supply Chain IT\\[3pt]
\textbf{Intelligence}    & Vision + OCR + LLM (cascade multi-niveaux)\\[3pt]
\textbf{Intégrations}    & Jira Cloud, Slack, Firebase, SendGrid\\[3pt]
\textbf{Date}            & \today
\end{tabular}

\vfill
{\color{primary!70}\small Document généré pour présentation mentor / soutenance}
\end{center}

\newpage

% ---------- Sommaire -------------------------------------------------
\tableofcontents
\newpage

% =====================================================================
\section{Introduction \& contexte}

\textbf{StockIT} est une application Android conçue pour digitaliser et
automatiser la gestion du parc informatique d'une équipe logistique.
Elle remplace les processus manuels (feuilles Excel, saisies répétées,
étiquetage papier) par un flux \emph{scan + IA} qui réduit
drastiquement le temps de traitement d'une réception ou d'une sortie
de matériel.

\begin{infobox}[Problème adressé]
Le service logistique reçoit chaque semaine des cartons de matériel
informatique (écrans, claviers, souris, laptops\dots). Chaque article
doit être~:
\begin{itemize}
  \item rattaché à un bon de commande (PO) et à une facture ;
  \item enregistré en stock avec toutes ses caractéristiques ;
  \item ensuite attribué à un ticket de support ouvert par un
        collaborateur.
\end{itemize}
Ces étapes, faites manuellement, sont \textbf{longues, sources
d'erreurs et non traçables}. StockIT automatise l'ensemble de la chaîne.
\end{infobox}

\subsection{Objectifs}
\begin{itemize}
  \item \textbf{Zéro double saisie} grâce à l'OCR et à la vision par IA.
  \item \textbf{Traçabilité totale} : chaque mouvement est journalisé
        et notifié en temps réel.
  \item \textbf{Assistance IA} pour identifier les produits, lire les
        factures, choisir le bon ticket et répondre aux questions.
  \item \textbf{Robustesse hors-ligne} : les fonctions critiques
        continuent de fonctionner sans connexion.
\end{itemize}

% =====================================================================
\section{Architecture générale}

\subsection{Vue en couches}

\begin{center}
\begin{tikzpicture}[node distance=4mm]
  \node[layer, fill=primary!15] (ui) {\textbf{Présentation}\\Activities, Adapters, Widget};
  \node[layer, below=of ui]      (ctl){\textbf{Contrôleurs}\\MainController, Coordinators, Workers};
  \node[layer, below=of ctl]     (mdl){\textbf{Modèle}\\Room (SQLite) — Entités \& DAOs};
  \node[layer, below=of mdl, fill=red!8, draw=accent] (svc){\textbf{Services externes}\\Claude, Gemini, Jira, Slack, Firebase, SendGrid};
  \foreach \a/\b in {ui/ctl,ctl/mdl,mdl/svc}
    \draw[arr] (\a) -- (\b);
\end{tikzpicture}
\end{center}

\subsection{Chiffres clés du projet}

\begin{center}
\begin{tabular}{lr}
\toprule
\textbf{Élément} & \textbf{Volume} \\
\midrule
Écrans (Activities Android)              & \textbf{15} \\
Adaptateurs \& contrôleurs UI            & \textbf{18} \\
Entités persistées (Room)                & \textbf{13} \\
Services externes intégrés               & \textbf{9}  \\
Niveaux de fallback dans la cascade IA   & \textbf{5}  \\
Version du schéma de base                & \textbf{22} \\
\bottomrule
\end{tabular}
\end{center}

% =====================================================================
\section{Panorama des fonctionnalités}

\subsection{Écrans principaux}

\begin{tabularx}{\linewidth}{@{}lX@{}}
\toprule
\textbf{Écran} & \textbf{Rôle}\\
\midrule
\textit{Login}                & Connexion par mot de passe ou empreinte biométrique, avec 3 rôles (Admin, Manager, Viewer). \\
\textit{Dashboard}            & Tableau de bord central : stock, achats, rapports, audit, utilisateurs, IA, classement. \\
\textit{Scan produit}         & Prise de photo + reconnaissance automatique du type d'équipement. \\
\textit{Scan lot (batch)}     & Détection de plusieurs codes-barres en continu à travers la caméra. \\
\textit{Réception colis}      & Photo du bon de livraison / facture avec extraction structurée. \\
\textit{Sélection PO}         & Choix du bon de commande parmi ceux détectés sur la facture. \\
\textit{Étiquette carton}     & Lecture automatique de la référence, quantité, marque, PO imprimé. \\
\textit{Scan sortie}          & Sortie de matériel avec choix du mode d'attribution à un ticket. \\
\textit{Liste tickets}        & Recherche et sélection d'un ticket Jira ouvert. \\
\textit{Assistant IA}         & Chat conversationnel en langage naturel sur l'état du stock. \\
\textit{Réclamations}         & Suivi des tickets internes de réclamation (ouvert / en cours / résolu). \\
\textit{Profil}               & Édition du profil, points, badges et historique. \\
\textit{Widget d'accueil}     & Vignette Android affichant le nombre d'articles en stock, mise à jour automatiquement. \\
\bottomrule
\end{tabularx}

\subsection{Fonctionnalités groupées par thème}

\begin{center}
\begin{tabularx}{\linewidth}{@{}lX@{}}
\toprule
\textbf{Thème} & \textbf{Fonctionnalités}\\
\midrule
\textbf{Réception}       & Scan facture, extraction PO/fournisseur, cross-check étiquette, entrée stock automatique.\\
\textbf{Sortie}          & Scan produit, saisie quantité, 3 modes d'attribution ticket, décrément stock.\\
\textbf{Tickets Jira}    & Lecture, filtrage, sélection, création automatique en cas d'anomalie.\\
\textbf{Assistant IA}    & Q\&R en français sur le stock, contexte inventaire injecté.\\
\textbf{Réclamations}    & Formulaire priorisé (Basse / Moyenne / Haute) + suivi statut.\\
\textbf{Utilisateurs}    & Rôles, authentification biométrique, gamification (points, badges, quêtes).\\
\textbf{Dashboard}       & KPIs, graphiques de tendance, articles sous seuil, historiques.\\
\textbf{Widget}          & Vignette temps réel sur l'écran d'accueil Android.\\
\textbf{Notifications}   & Push FCM + alertes Slack sur chaque mouvement critique.\\
\textbf{Rapports}        & PDF mensuel généré et envoyé automatiquement par email.\\
\textbf{Scan en lot}     & Détection en continu de plusieurs codes-barres via la caméra.\\
\textbf{Audit \& logs}   & Chaque action journalisée avec utilisateur, date, appareil.\\
\textbf{Gamification}    & Quêtes, points, badges, classement des utilisateurs.\\
\bottomrule
\end{tabularx}
\end{center}

% =====================================================================
\section{Workflow \textsc{Entrée} — réception d'un article}

\begin{infobox}[Objectif]
Identifier l'équipement reçu, le rattacher à son bon de commande et à
sa facture, puis l'enregistrer en stock — le tout en quelques photos.
\end{infobox}

\begin{enumerate}
  \item \textbf{Scan du produit.} Le technicien prend une photo de
        l'article. L'IA reconnaît automatiquement le type d'équipement
        (écran, clavier, souris\dots).
  \item \textbf{Scan de la facture / bon de livraison.} Une seconde
        photo permet d'extraire, par reconnaissance de texte, le
        fournisseur et la liste des commandes présentes sur le document.
  \item \textbf{Sélection de la commande.} L'application affiche les
        commandes détectées et met en évidence celle qui correspond au
        produit scanné (\emph{En base}, \emph{Correspond au scan}).
  \item \textbf{Scan de l'étiquette du carton.} La quantité, la
        référence article et la marque sont lues automatiquement, puis
        comparées à la commande sélectionnée.
  \item \textbf{Enregistrement en stock.} L'entrée est validée : le
        stock est mis à jour, une trace est conservée dans l'historique
        et une notification est envoyée à l'équipe.
\end{enumerate}

\begin{center}
\begin{tikzpicture}[node distance=5mm and 6mm]
  \node[box] (n1) {Scan\\produit};
  \node[box, right=of n1] (n2) {Scan\\facture};
  \node[box, right=of n2] (n3) {Choix de\\la commande};
  \node[box, right=of n3] (n4) {Scan\\étiquette};
  \node[box, right=of n4] (n5) {Entrée\\en stock};
  \foreach \a/\b in {n1/n2,n2/n3,n3/n4,n4/n5}
    \draw[arr] (\a) -- (\b);
\end{tikzpicture}
\end{center}

% =====================================================================
\section{Workflow \textsc{Sortie} — attribution à un ticket}

\begin{infobox}[Objectif]
Sortir un ou plusieurs articles du stock et tracer chaque unité vers
son destinataire, via l'un des \textbf{trois modes d'attribution}.
\end{infobox}

\begin{enumerate}
  \item \textbf{Scan du produit à sortir} : l'IA identifie l'équipement.
  \item \textbf{Saisie de la quantité} : contrôlée contre le stock disponible.
  \item \textbf{Choix du mode d'attribution} : pour chaque unité,
        l'utilisateur choisit l'un des trois modes ci-dessous.
\end{enumerate}

\vspace{4pt}
\begin{center}
\begin{tikzpicture}[node distance=5mm and 8mm]
  \node[box] (s1) {Scan\\produit};
  \node[box, right=of s1] (s2) {Quantité\\à sortir};
  \node[cbox, right=of s2] (s3) {Choix\\du mode};
  \node[cbox, right=13mm of s3, yshift=12mm] (m1) {Mode 1\\Automatique};
  \node[cbox, right=13mm of s3]              (m2) {Mode 2\\Liste};
  \node[cbox, right=13mm of s3, yshift=-12mm] (m3) {Mode 3\\Identifiant};
  \node[box,  right=10mm of m2]              (out) {Sortie\\enregistrée};
  \draw[arr] (s1) -- (s2);
  \draw[arr] (s2) -- (s3);
  \draw[arr] (s3) -- (m1);
  \draw[arr] (s3) -- (m2);
  \draw[arr] (s3) -- (m3);
  \draw[arr] (m1.east) -| (out.north);
  \draw[arr] (m2)      -- (out);
  \draw[arr] (m3.east) -| (out.south);
\end{tikzpicture}
\end{center}

\subsection{Les 3 modes d'attribution}

\begin{center}
\begin{tabularx}{\linewidth}{@{}lX@{}}
\toprule
\textbf{Mode}            & \textbf{Description}\\
\midrule
\textbf{1. Automatique (IA)} & L'application propose elle-même le
    ticket le plus pertinent en analysant les demandes ouvertes
    (résumé, priorité, échéance) et en les faisant correspondre au
    produit scanné.\\
\addlinespace
\textbf{2. Liste}            & Le technicien parcourt la liste des
    tickets ouverts, filtre par mot-clé ou priorité, et sélectionne
    manuellement celui qui correspond.\\
\addlinespace
\textbf{3. Identifiant}      & Le technicien tape directement le
    numéro d'un ticket connu ; l'application vérifie qu'il existe
    avant de valider.\\
\bottomrule
\end{tabularx}
\end{center}

\begin{enumerate}\setcounter{enumi}{3}
  \item \textbf{Validation} : la sortie est enregistrée avec le lien
        vers le ticket choisi et la boucle recommence tant qu'il reste
        des unités à attribuer.
  \item \textbf{Clôture} : récapitulatif à l'écran + notification
        automatique de l'équipe.
\end{enumerate}

% =====================================================================
\section{La cascade IA — jamais bloquante}

Chaque appel de reconnaissance passe par 5 niveaux successifs. Dès
qu'un niveau répond, les suivants ne sont pas sollicités. Le dernier
niveau est \textbf{100\,\% local}, garantissant que l'application
fonctionne même sans réseau.

\begin{center}
\begin{tikzpicture}[node distance=3mm]
  \node[layer] (l1) {1. Passerelle Cimpress — Claude Opus~4};
  \node[layer, below=of l1] (l2) {2. Portkey (routeur multi-LLM)};
  \node[layer, below=of l2] (l3) {3. Google Gemini};
  \node[layer, below=of l3] (l4) {4. HuggingFace (\emph{vit-base})};
  \node[layer, below=of l4, fill=green!8, draw=green1]
       (l5) {5. ML~Kit local — \emph{offline, jamais bloquant}};
  \foreach \a/\b in {l1/l2,l2/l3,l3/l4,l4/l5}
    \draw[arr] (\a) -- (\b);
\end{tikzpicture}
\end{center}

\begin{ideabox}[Pourquoi cette cascade ?]
\begin{itemize}
  \item \textbf{Qualité} maximale en tête de chaîne (LLM vision de
        dernière génération).
  \item \textbf{Résilience} : coupure réseau, quota, erreur API — un
        niveau prend toujours le relais.
  \item \textbf{Coût maîtrisé} : les niveaux gratuits (ML~Kit) sont
        des filets de sécurité, pas la solution principale.
\end{itemize}
\end{ideabox}

% =====================================================================
\section{Automatisations \& scénarios intelligents}

\subsection{Kit anti-oubli}
Lorsqu'un \textbf{écran} est reçu, un dialogue force la vérification
des accessoires obligatoires : \emph{câble d'alimentation} et
\emph{câble HDMI}. Si l'un manque, une alerte est envoyée dans Slack
et un ticket Jira d'anomalie est créé automatiquement — sans bloquer
l'entrée en stock.

\subsection{Auto-matching ticket $\leftrightarrow$ produit}
En mode automatique de sortie, l'IA reçoit la liste des tickets
ouverts (résumé, priorité, échéance) et l'équipement scanné, puis
retourne une répartition des unités en tenant compte de la priorité
et de la date d'échéance.

\subsection{Création automatique de tickets Jira}
Trois déclencheurs :
\begin{itemize}
  \item Kit accessoire incomplet (écran sans câbles).
  \item Stock d'un article sous seuil critique.
  \item Bon de commande non reçu au-delà d'un délai.
\end{itemize}

\subsection{Notifications Slack contextuelles}
Chaque action génère un message ciblé :
\begin{itemize}
  \item Succès : \emph{« \checkmark~Sortie enregistrée — 3 claviers vers TICKET-42 »}
  \item Anomalie : \emph{« \faExclamationTriangle~Kit incomplet — écran sans câble HDMI »}
  \item Rapport : \emph{« \faFilePdf~Rapport mensuel prêt, 145 articles »}
\end{itemize}

\subsection{Rapports PDF planifiés}
Un travail de fond (WorkManager) s'exécute à intervalle
configurable : il collecte l'état du stock, génère un PDF de
synthèse (avec résumé rédigé par l'IA) et l'envoie par email via
SendGrid — même si l'application est fermée.

\subsection{Optimisation d'image avant l'IA}
Les photos sont compressées à 1024\,px et 85\,\% de qualité avant
envoi aux modèles de vision — latence réduite et consommation
réseau maîtrisée.

\subsection{Cohérence PO facture $\leftrightarrow$ étiquette}
Le PO imprimé sur l'étiquette du carton est comparé au PO sélectionné
sur la facture. Un avertissement non-bloquant est affiché en cas
d'incohérence.

\subsection{Widget synchronisé en temps réel}
Chaque mouvement de stock déclenche un rafraîchissement du widget
d'accueil : la vignette affiche toujours le compteur exact.

\subsection{Gamification}
Chaque scan et chaque action utilisateur rapportent des points, des
badges et alimentent un système de quêtes et de classement.

% =====================================================================
\section{Modèle de données}

\begin{center}
\begin{tabularx}{\linewidth}{@{}lX@{}}
\toprule
\textbf{Entité} & \textbf{Rôle}\\
\midrule
\textit{Product}          & Article en stock (nom, catégorie, quantité, seuil, marque\dots)\\
\textit{PurchaseOrder}    & Bon de commande (fournisseur, quantité, statut)\\
\textit{ShippingOrder}    & Expédition sortante \\
\textit{StockMovement}    & Entrée / sortie (produit, quantité, ticket associé, motif)\\
\textit{Category}         & Catégorie de produit \\
\textit{Supplier}         & Fournisseur \\
\textit{User}             & Utilisateur (rôle, points, badges, niveau)\\
\textit{AuditLog}         & Journal d'audit de toutes les actions \\
\textit{Claim}            & Réclamation (sujet, priorité, statut)\\
\textit{Quest}            & Quête de gamification (objectif, progression, récompense)\\
\textit{AnalyzedTicket}   & Cache local d'un ticket Jira analysé \\
\textit{ChatMessage}      & Message de l'assistant IA \\
\textit{AIInsight}        & Recommandation ou anomalie détectée par l'IA \\
\bottomrule
\end{tabularx}
\end{center}

% =====================================================================
\section{Intégrations \& services externes}

\begin{center}
\begin{tabularx}{\linewidth}{@{}lXl@{}}
\toprule
\textbf{Service} & \textbf{Usage} & \textbf{Auth}\\
\midrule
Passerelle Cimpress (Claude Opus~4) & Vision produit, extraction facture, matching ticket & Bearer token \\
Portkey                             & Routeur LLM secondaire (chat, fallback) & API key \\
Google Gemini                       & Vision de secours & API key \\
HuggingFace (\emph{vit-base})       & Classification d'image de secours & Bearer token \\
ML~Kit (Google, local)              & OCR, codes-barres, image labeling hors-ligne & --- \\
Jira Cloud REST v3                  & Lecture et création de tickets & Basic (email + token) \\
Firebase Cloud Messaging            & Notifications push temps réel & Compte de service \\
Slack \emph{Incoming Webhooks}      & Alertes automatiques équipe & Webhook URL \\
SendGrid                            & Envoi des rapports PDF par email & API key \\
\bottomrule
\end{tabularx}
\end{center}

% =====================================================================
\section{Outils \& bibliothèques utilisés}

\begin{tabularx}{\linewidth}{@{}lX@{}}
\toprule
\textbf{Catégorie} & \textbf{Technologies}\\
\midrule
\textbf{Langage \& build}      & Java 11, Gradle Kotlin DSL, Android Gradle Plugin, min~SDK~26, target~SDK~35 \\
\textbf{UI}                    & Material Design, AppCompat, RecyclerView, ConstraintLayout, CardView, DrawerLayout, MPAndroidChart (graphiques)\\
\textbf{Caméra \& vision}      & CameraX (core, camera2, lifecycle, view), ML~Kit (Text Recognition, Barcode Scanning, Image Labeling, Object Detection)\\
\textbf{IA / LLM}              & Client OkHttp vers Claude, Gemini, Portkey, HuggingFace ; SDK Google Generative AI\\
\textbf{Base de données}       & Room (SQLite), 22 versions de schéma \\
\textbf{Réseau}                & OkHttp, Retrofit, Gson\\
\textbf{Backend cloud}         & Firebase BoM (Messaging, Firestore optionnel)\\
\textbf{PDF \& rapports}       & iText 7 (\emph{itext7-core}) \\
\textbf{Email}                 & SendGrid Java SDK \\
\textbf{Tâches d'arrière-plan} & WorkManager (PeriodicWorkRequest) \\
\textbf{Widget}                & AppWidgetProvider (framework Android) \\
\textbf{Authentification}      & AndroidX Biometric (empreinte / visage) \\
\textbf{Tests}                 & JUnit 4, AndroidX Test, Espresso \\
\bottomrule
\end{tabularx}

% =====================================================================
\section{Sécurité \& bonnes pratiques}

\subsection{Gestion des secrets}
Les clés d'API sont injectées dans \emph{BuildConfig} à la
compilation, à partir de variables d'environnement (fichiers
\emph{set\_env.ps1} / \emph{.env.sh} non versionnés). Aucun secret
n'est écrit en clair dans le dépôt Git.

\subsection{Authentification}
Deux modes disponibles : mot de passe et empreinte biométrique (via
l'API AndroidX Biometric). Trois rôles applicatifs : \emph{Admin},
\emph{Manager}, \emph{Viewer}.

\subsection{Traçabilité}
Chaque action utilisateur est enregistrée dans une table d'audit
(utilisateur, action, date, appareil) — non modifiable depuis
l'interface. Cette table alimente également l'onglet \emph{Audit} du
dashboard.

\subsection{Permissions Android}
Camera, Internet, État réseau, Notifications, Biométrie — toutes
demandées à l'exécution et strictement limitées à l'usage nécessaire.

\subsection{Réseau}
Timeouts explicites (connexion 10\,s, lecture 30\,s+), HTTPS
systématique, gestion défensive des erreurs (aucun \emph{crash} sur
échec IA — la cascade prend le relais).

% =====================================================================
\section{Valeur ajoutée \& innovations}

\begin{ideabox}[Points différenciants]
\begin{enumerate}
  \item \textbf{Cascade IA à 5 niveaux, jamais bloquante} — l'app
        continue de fonctionner même hors ligne.
  \item \textbf{Zéro double saisie} — vision + OCR partout : facture,
        étiquette, produit.
  \item \textbf{Attribution intelligente ticket $\leftrightarrow$ matériel} —
        priorisation automatique par urgence et échéance.
  \item \textbf{Traçabilité totale} — audit + Slack + Jira synchronisés
        à chaque action.
  \item \textbf{Assistant IA embarqué} — questions en langage naturel
        sur l'état du stock.
  \item \textbf{Kit anti-oubli} — vérification obligatoire des
        accessoires (câbles) sur les écrans.
  \item \textbf{Rapports PDF planifiés} envoyés automatiquement par
        email (WorkManager + SendGrid).
  \item \textbf{Widget d'accueil} synchronisé après chaque mouvement
        de stock.
  \item \textbf{Gamification} — quêtes, points et badges pour
        engager les utilisateurs logistique.
\end{enumerate}
\end{ideabox}

% =====================================================================
\section{Conclusion}

StockIT démontre qu'une application mobile bien pensée peut, à elle
seule, remplacer un ensemble de processus manuels dans la gestion
d'un parc informatique. En combinant \textbf{vision par IA},
\textbf{OCR local}, \textbf{intégrations métier} (Jira, Slack) et
\textbf{automatisations d'arrière-plan}, elle fait gagner du temps à
l'équipe logistique tout en garantissant une traçabilité complète.

L'architecture modulaire (cascade IA remplaçable, entités Room
extensibles, services externes découplés) permet d'ajouter facilement
de nouvelles fonctionnalités : lecture de bons Zycus, prévision de
consommation, ou intégration ERP.

\vspace{6pt}
\begin{center}
\color{primary!70}\itshape\small
Fin du rapport — merci de votre attention.
\end{center}

\end{document}
